% Chellas's basic conditional logic CK in the theorem-only Hilbert form.
\calculussheet{sheet:hilbert-conditional-ck}{Chellas Basic Conditional Logic \(\mathsf{CK}\)}{Language: classical propositional formulas extended by the binary conditional \(A>B\); theoremhood is \(\derives_{\mathsf{CK}}A\); all substitutions are uniform and capture avoiding; the two conditional rules are theorem-only; equality: none.}
\noindent\emph{Controlling source:} \cite{Chellas1975}.
\hbrule{rule:p08-ck-cpl}{CPL base}{}{A}{\(A\) is an instance of a classical propositional tautology; modus ponens is inherited from the fixed Hilbert CPL base}
\hbrule{rule:p08-ck-rcea}{RCEA}{\derives_{\mathsf{CK}}A\leftrightarrow B}{\derives_{\mathsf{CK}}((A>C)\leftrightarrow(B>C))}{theorem-only replacement of equivalent antecedents}
\hbrule{rule:p08-ck-rck}{RCK}{\derives_{\mathsf{CK}}((B_1\land\cdots\land B_n)\to B)}{\derives_{\mathsf{CK}}(((A>B_1)\land\cdots\land(A>B_n))\to(A>B))}{\(n\geq0\); for \(n=0\), the premise is \(\derives B\) and the conclusion is \(\derives A>B\)}
